% =============================================================================
% literature_review.tex -- Literature Review (author: Yogesh Modi, 2401010519)
% References: docs/report/refs_yogesh.bib
% =============================================================================
\section{Literature Review}
\label{sec:literature}

\subsection{The Classic Five-Stage RISC Pipeline}
Reduced Instruction Set Computer (RISC) design was shaped by early academic
processors such as Berkeley RISC~I, which showed that a small set of simple,
fixed-length instructions with a load/store architecture and a large register
file can be implemented efficiently in VLSI~\cite{ym_patterson_risc1}. Simple,
regular instructions make pipelining natural. The textbook organisation used in
teaching and in many small cores is the five-stage pipeline of Patterson and
Hennessy~\cite{ym_patterson_cod_riscv}: Instruction Fetch (IF), Instruction
Decode and register read (ID), Execute or address calculation (EX), Memory
access (MEM) and Write Back (WB), separated by the pipeline registers IF/ID,
ID/EX, EX/MEM and MEM/WB. Ideally one instruction completes every clock cycle,
while each stage only has to do a fraction of the work of a single-cycle
design. Harris and Harris~\cite{ym_harris_ddca_riscv} present the same
organisation at the hardware-description level and build a pipelined RISC-V
processor step by step in SystemVerilog/VHDL, which makes it a direct guide for
an HDL implementation.

\subsection{The RISC-V Instruction Set Architecture}
RISC-V is an open, royalty-free instruction set architecture developed at the
University of California, Berkeley~\cite{ym_waterman_thesis}. Its design goals
include a small base integer ISA that is sufficient for a complete software
stack, optional standard extensions, and an encoding that simplifies hardware.
The unprivileged specification~\cite{ym_riscv_spec} defines the 32-bit base
integer set RV32I with 32 general-purpose registers, where register \texttt{x0}
is hard-wired to zero. Instructions use six formats (R, I, S, B, U and J).
The source and destination register fields are kept in the same bit positions in
every format, and the immediate sign bit is always instruction bit~31. This lets
the decoder read the register file in parallel with decoding the opcode and makes
the immediate generator a small multiplexer. These properties are exploited
directly by the Instruction Decode stage of our design.

\subsection{Pipeline Hazards and Their Handling}
Overlapping instructions introduces three kinds of
hazards~\cite{ym_patterson_cod_riscv}. \emph{Structural hazards} occur when two
instructions need the same hardware resource at once; they are avoided by
separate instruction and data memories (a Harvard organisation) and a register
file with independent read and write ports. \emph{Data hazards} occur when an
instruction needs a result that an earlier instruction has not yet written back.
The simplest fix is to \emph{stall} the pipeline by inserting bubbles, but
\emph{forwarding} (bypassing) the result from the EX/MEM or MEM/WB registers to
the ALU inputs removes most stalls. A load followed immediately by a dependent
instruction (the load-use case) still needs a one-cycle stall, because the
loaded value is only available after MEM. Writing the register file in the first
half of the cycle and reading it in the second half, or an equivalent
write-first bypass, removes the WB-to-ID case. \emph{Control hazards} arise from
branches and jumps, whose outcome is not known when the next instruction is
fetched. Options include stalling, a static predict-not-taken policy that
flushes the wrongly fetched instructions when a branch is taken, and dynamic
branch prediction. Smith's study of branch prediction
strategies~\cite{ym_smith_branch} showed that simple history-based predictors
achieve high accuracy and motivated the counter-based predictors used in modern
cores.

\subsection{RISC-V Processor Implementations}
Several open RISC-V cores show how pipeline depth is traded against area, power
and performance. Gautschi \emph{et al.}~\cite{ym_gautschi_ri5cy} present RI5CY, a
four-stage in-order RISC-V core with digital-signal-processing extensions such as
hardware loops and SIMD operations, designed for near-threshold operation in
energy-constrained Internet-of-Things (IoT) devices. Schiavone \emph{et
al.}~\cite{ym_schiavone_ulp} compare RI5CY with smaller cores such as the
two-stage Zero-riscy and Micro-riscy, and show that fewer pipeline stages and
fewer features reduce area and power at the cost of performance. At the
higher-performance end, the Rocket Chip generator~\cite{ym_asanovic_rocket}
produces a five-stage in-order RISC-V core and complete SoCs from a
parameterised Chisel description, and the Berkeley Out-of-Order Machine
(BOOM)~\cite{ym_celio_boom} is a synthesizable, parameterised out-of-order
RISC-V core. These works show the breadth of microarchitectures that the open
ISA supports. They are, however, large, highly optimised designs that are not
intended as a first study of pipelining.

The open-source project ProtocolX~\cite{ym_protocolx_repo}, which our team
studied as a reference, implements a five-stage RV32I pipeline in Verilog HDL
with a forwarding unit for read-after-write hazards, a one-cycle stall for
load-use hazards and a flush of wrongly fetched instructions on taken branches,
and it is simulated with Icarus Verilog. Dynamic branch prediction is listed
there only as future work.

\subsection{Relation to Our Work}
Our processor follows the classic five-stage organisation of
\cite{ym_patterson_cod_riscv,ym_harris_ddca_riscv} and the RV32I encoding
of~\cite{ym_riscv_spec}. Unlike the industrial cores
of~\cite{ym_gautschi_ri5cy,ym_schiavone_ulp,ym_asanovic_rocket,ym_celio_boom},
its scope is deliberately limited to a clear, fully verified teaching
implementation: a defined RV32I subset (R-type and I-type ALU operations, LW,
SW, BEQ, BNE, BLT, BGE, JAL, JALR and LUI), EX/MEM and MEM/WB forwarding, a
one-cycle load-use stall, a write-first register-file bypass and a static
predict-not-taken branch policy with a two-cycle flush. Compared with the
reference project~\cite{ym_protocolx_repo}, our design is written from scratch
against a shared interface specification, divided into independently owned
modules, and every module has its own self-checking unit testbench in addition
to the system-level tests. Dynamic branch prediction, the M extension,
interrupts and FPGA synthesis are outside the Phase~I scope.
